In 1991 schreef een Finse student met de naam Linus Torvalds\index{Torvalds, Linus} een berichtje in een news-group dat hij bezig was met een operating system waarvoor hij de kernel al had geschreven en de GNU-tools gebruikte om er een werkbaar systeem van te maken. Linus was student en wist dat de kosten voor een besturingssysteem voor studenten soms moeilijk op te brengen was. Hij wilde dat zijn systeem gratis zou zijn.

De kernel die Linus geschreven had ging naar de maker heten en omdat alle Unix-afgeleiden op een X eindigen werd dat dus Linux.

Linux is zoals gezegd een kernel; een abstractie laag tussen de hardware en de gebruikers interface. Voor een compleet systeem is er veel meer nodig. De overige software op wat wij nu een Linux systeem noemen is dan ook veelal afkomstig van het GNU-project. Sommige mensen willen dan ook graag dat je spreekt van GNU/Linux, maar de hele wereld spreekt meestal gewoon over Linux als we een systeem bedoelen met een Linux kernel.

